\chapter{Statistical Analysis}
\label{chap:statistical-analysis}

\section{Overview of the Statistical Design}
\label{sec:statistical-design-overview}

The analysis was designed as an exploratory
full-dataset statistical association analysis, to test whether network properties differed systematically between
the interictal baseline window and the preictal windows T0, T1, and T2.

Each graph-level observation was derived from one functional connectivity
network. A network was uniquely identified by the subject/case identifier,
seizure identifier, temporal window, frequency band, connectivity method, and
graph metric. The statistical analysis preserved this structure so that graph
metrics were not treated as independent biological samples simply because they
were produced across multiple frequency bands, connectivity methods, or
network measures.

The main experimental factor was temporal window, with four levels:
Baseline, T0, T1, and T2. Baseline represented an interictal reference window,
whereas T0, T1, and T2 represented consecutive preictal windows approaching
seizure onset. The statistical question was whether graph metrics varied across
these windows, and in particular whether the preictal windows differed from
Baseline.

The primary analysis was performed at the seizure level. In this analysis,
each seizure was treated as the repeated-measures unit because the same seizure
contributed graph metrics across multiple temporal windows. Seizures were
nested within subject/case identifiers to account for the fact that multiple
seizures could come from the same CHB-MIT case folder. This design helped avoid
treating seizures from the same subject/case identifier as fully independent
observations.

In addition to the seizure-level analysis, a subject-average sensitivity
analysis was performed. In this analysis, graph metrics were first averaged
across seizures within each subject/case identifier for each temporal window,
frequency band, connectivity method, and graph metric. This produced one
average trajectory per subject/case identifier. The purpose of this analysis
was to evaluate whether effects observed at the seizure level were also visible
when each subject/case identifier contributed equally to the group-level
analysis.

For each graph metric, connectivity method, and frequency band, three types of
statistical tests were performed. First, an omnibus window-effect test assessed
whether there was evidence of any difference across Baseline, T0, T1, and T2.
Second, planned post-hoc paired comparisons tested Baseline against each
preictal window. Third, a preictal trend analysis examined whether graph
metrics changed monotonically across T0, T1, and T2 as seizure onset
approached.

Overall, the statistical design combined a primary seizure-level analysis with
a subject-average sensitivity analysis. This allowed the thesis to examine both
seizure-level temporal dynamics and the robustness of these effects at the
subject/case level, while maintaining a clear distinction between exploratory
statistical association and predictive modelling.

\section{Data Organization for Statistical Testing}
\label{sec:data-organization-statistical-testing}

The statistical analysis used the graph feature outputs described in
\cref{sec:connectivity-graph-pipeline-output}. Each row of the original graph
feature table represented one functional network derived from a specific
combination of subject/case identifier, seizure, temporal window, frequency
band, and connectivity method. The graph metrics were then converted into a
long-format table, where each row contained one value for one graph metric.

The long-format organization made it possible to analyze each graph metric,
connectivity method, and frequency band combination separately. For every such
combination, the statistical model compared values across temporal windows.
This avoided combining graph metrics with different meanings, connectivity
methods with different physiological interpretations, or frequency bands with
different spectral content.

The identifiers retained for each statistical row were: subject/case identifier,
seizure identifier, temporal window, frequency band, connectivity method, graph metric. graph metric value.

This organization preserved the hierarchical structure of the dataset.
Observations were grouped by seizure and subject/case identifier rather than
treated as unrelated measurements.

\subsection{Seizure-Level Analysis}
\label{subsec:seizure-level-analysis}

The seizure-level analysis was the primary statistical analysis. In this
analysis, each seizure contributed repeated measurements across temporal
windows. For a given graph metric, connectivity method, and frequency band, the
same seizure could contribute values for Baseline, T0, T1, and T2. This made it
possible to test whether the graph metric changed within seizures as the EEG
moved from the interictal baseline window toward seizure onset.

In the primary model, seizures were treated as repeated-measures units nested
within subject/case identifiers. This was important because several CHB-MIT
case folders contain multiple seizures. Treating all seizure-level rows as
fully independent would overstate the amount of independent biological
information in the dataset. The mixed-model structure was therefore used to
account for repeated seizure observations while retaining the seizure-level
temporal information.

The seizure-level analysis is useful because it preserves variability between
seizures within the same subject/case identifier. This is relevant for epilepsy
data, where different seizures from the same individual may show different
preictal dynamics. However, because some subject/case identifiers contributed
more seizures than others, a subject-average sensitivity analysis was also
performed.

\subsection{Subject-Average Sensitivity Analysis}
\label{subsec:subject-average-sensitivity-analysis}

The subject-average sensitivity analysis was performed to evaluate whether
patterns observed at the seizure level were also visible when each subject/case
identifier contributed equally. In this analysis, graph metric values were
first averaged across seizures within each subject/case identifier for each
temporal window, frequency band, connectivity method, and graph metric.

The resulting subject-average table contained one average value per
subject/case identifier and temporal window for each metric--method--band
combination. Standard deviation across seizures and the number of contributing
seizures were also retained as descriptive information. The statistical tests
were then repeated using these subject-level average trajectories.

This analysis reduced the influence of subject/case identifiers with many
seizures. For example, a subject/case identifier contributing many usable
seizures could strongly influence a seizure-level analysis, whereas in the
subject-average analysis it contributed one average trajectory, like every
other subject/case identifier. The subject-average approach therefore provided
a complementary check on whether the main findings reflected consistent
subject-level patterns rather than being driven mainly by subjects with larger
numbers of seizures.

Averaging across seizures reduces
within-subject seizure variability and may hide seizure-specific preictal
effects. For this reason, the thesis reports the seizure-level analysis as the
primary analysis and uses the subject-average analysis to assess robustness at
the subject/case level.

\section{Outcome Variables and Experimental Factors}
\label{sec:outcome-variables-experimental-factors}

The statistical analysis was organized around graph-level outcome variables and
a fixed set of experimental factors. Each outcome variable represented one
graph-theoretical feature extracted from a functional connectivity network.
The experimental factors described how each network was generated: the temporal
window, frequency band, and connectivity method.

The outcome variables analyzed in this thesis were the graph metrics described
in \cref{sec:graph-theoretical-metrics}.

Although graph density was computed and retained as a diagnostic variable, it
was not treated as a primary statistical outcome. This is because the graph
construction procedure applied a proportional threshold targeting a fixed
density. Therefore, density mainly served as a quality-control measure
confirming the realized number of retained edges after thresholding.

The statistical tests therefore evaluated whether graph metrics changed from
Baseline to the preictal period and whether there was evidence of progressive
change within the preictal interval.

The second experimental dimension was frequency band. The analysis was
performed separately for the six frequency bands used throughout the pipeline:
delta, theta, alpha1, alpha2, beta1, and beta2. This band structure follows the
pre-seizure graph-theoretical analysis of Vecchio et al., who analyzed
connectivity separately across the same frequency ranges \cite{Vecchio2016}.
The bands were not pooled in the statistical models because seizure-related
network organization may differ across frequency ranges, and combining bands
could obscure frequency-specific effects \cite{Ponten2007,Lanzone2022}.

The third experimental dimension was connectivity method. The analysis was
performed separately for coherence, wPLI, and AEC networks.
Separate statistical testing by method allowed to evaluate whether
observed window effects were specific to a particular definition of functional
connectivity.

For each combination of graph metric, connectivity method, and frequency band,
a separate statistical analysis was performed. Since the final analysis
included \num{8} graph metrics, \num{3} connectivity methods, and \num{6}
frequency bands, this produced \num{144} metric--method--band combinations. Each combination was analyzed independently
across temporal windows. This design made it possible to identify which graph
properties, connectivity definitions, and frequency bands showed evidence of
preictal change.

The statistical models did not treat individual epochs, channels, graph
metrics, frequency bands, or connectivity methods as independent biological
samples. Instead, the unit of repeated measurement was defined at the seizure
level for the primary analysis and at the subject/case level for the
sensitivity analysis. This distinction was important because the large number
of derived graph features does not imply an equally large number of independent
subjects or seizures.

\section{Omnibus Testing Across Temporal Windows}
\label{sec:omnibus-testing-temporal-windows}

The first statistical question was whether each graph metric showed an overall
effect of temporal window. This was tested using an omnibus window-effect test
for each graph metric, connectivity method, and frequency band combination.

The purpose of the omnibus test was to determine whether there was evidence of
any difference across the four temporal windows, without first specifying which
preictal window differed from Baseline. This step was important because the
analysis involved many graph metrics, connectivity methods, and frequency
bands. Planned pairwise comparisons were therefore interpreted mainly when the
corresponding omnibus test indicated a significant window effect after
multiple-comparison correction.

For the primary seizure-level analysis, the omnibus test was performed using a
linear mixed-effects model. This modelling framework was selected because it is
appropriate for repeated-measures and grouped data, allowing temporal-window
effects to be estimated while accounting for the non-independence of seizures
within subject/case identifiers \cite{Laird1982,Pinheiro2000}. For each metric--method--band
combination, temporal window was included as a categorical fixed effect:

\[
    y_{s,w}
    =
    \beta_0
    +
    \beta_w
    +
    u_{\mathrm{subject}}
    +
    u_{\mathrm{seizure}}
    +
    \varepsilon_{s,w},
\]

where $y_{s,w}$ is the graph metric value for seizure $s$ in window $w$,
$\beta_0$ is the intercept, $\beta_w$ represents the fixed effect of temporal
window, $u_{\mathrm{subject}}$ accounts for grouping by subject/case identifier,
$u_{\mathrm{seizure}}$ accounts for repeated measurements within seizure, and
$\varepsilon_{s,w}$ is the residual error term.

Baseline was used as the reference window in the mixed model. The omnibus
window effect was then evaluated using a Wald chi-square test of the window
terms \cite{Wald1943}. This tested whether the temporal window coefficients were jointly
different from zero. A significant result indicated that the graph metric
varied across Baseline, T0, T1, and T2 for the corresponding connectivity
method and frequency band.

When the mixed-effects model could not be estimated, the pipeline attempted a
repeated-measures fallback analysis. If the Pingouin statistical package was
available, a one-way repeated-measures ANOVA was attempted; otherwise, a manual
one-way repeated-measures ANOVA was available as a secondary fallback
\cite{Girden1992,Vallat2018}. This fallback preserved the within-unit
repeated-measures structure, although it did not explicitly model the nesting
of seizures within subject/case identifiers. Cases in which neither the
mixed-effects model nor the fallback analysis produced a valid test statistic
were marked as non-estimable and excluded from inferential interpretation.

The same omnibus testing logic was also applied to the subject-average
sensitivity analysis. In that case, the repeated-measures unit was the
subject/case identifier rather than the seizure. Each subject/case identifier
contributed one average value per window for a given metric--method--band
combination. The subject-average omnibus model therefore tested whether the
average subject-level trajectory differed across Baseline, T0, T1, and T2.

For each omnibus test, the pipeline recorded the number of complete repeated
observations, the number of contributing subject/case identifiers, the model
type, the test statistic, degrees of freedom where applicable, the uncorrected
$p$-value, the FDR-corrected $p$-value, and any fallback or skip reason. Tests
were skipped when the number of complete repeated observations was below the
predefined minimum required for analysis.

\section{Post-hoc Baseline-vs-Preictal Comparisons}
\label{sec:posthoc-baseline-preictal-comparisons}

After the omnibus temporal-window tests, planned post-hoc comparisons were
performed to identify which preictal windows differed from the interictal
baseline condition. These comparisons were defined a priori because the central
research question of the thesis concerned whether graph-theoretical network
features changed before seizure onset relative to Baseline.

For each graph metric, connectivity method, and frequency band combination,
three paired Baseline-vs-preictal contrasts were tested: Baseline vs. T0,
Baseline vs. T1,
Baseline vs. T2.

These contrasts compared the interictal reference window with each of the
three consecutive preictal windows. T0 represented the earliest preictal
window, T1 the middle preictal window, and T2 the window closest to seizure
onset. This design allowed the analysis to evaluate whether a metric differed
from Baseline at any specific preictal interval.

For each paired contrast, the primary post-hoc test was a paired-samples
$t$-test. If $x_{\mathrm{Baseline},s}$ denotes the baseline value for seizure
$s$ and $x_{\mathrm{Preictal},s}$ denotes the corresponding preictal value, the
paired difference was defined as

\[
    d_s =
    x_{\mathrm{Preictal},s}
    -
    x_{\mathrm{Baseline},s}.
\]

The paired $t$-test evaluated whether the mean of these within-seizure
differences differed from zero. A positive mean difference indicated that the
metric was higher in the preictal window than in Baseline, whereas a negative
mean difference indicated that the metric was lower in the preictal window.

For each contrast, the analysis recorded the number of paired observations, the
number of contributing subject/case identifiers, the Baseline mean, the
comparison-window mean, the mean difference, the $t$ statistic, degrees of
freedom, the uncorrected $p$-value, and the FDR-corrected $p$-value. 
Cohen's $d_z$ was also computed as an effect-size measure for paired data
\cite{Lakens2013}:

\[
    d_z =
    \frac{\bar{d}}{s_d},
\]

where $\bar{d}$ is the mean paired difference and $s_d$ is the standard
deviation of the paired differences.

As a non-parametric sensitivity check, a Wilcoxon signed-rank test was also
computed for each paired contrast when possible \cite{Wilcoxon1945}. The
Wilcoxon signed-rank test does not require the paired differences to be normally
distributed, although its standard interpretation assumes that the distribution
of paired differences is approximately symmetric. Therefore, it provided an
additional check on whether the paired comparison was robust to the normality
assumption of the paired $t$-test. Both the paired $t$-test and Wilcoxon test
results were retained in the statistical output. 
Both the paired $t$-test and Wilcoxon test results were retained
in the statistical output.

The same planned post-hoc contrasts were repeated for the subject-average
sensitivity analysis. In that analysis, the paired observations were
subject/case-level average values rather than individual seizure-level values.
Thus, each subject/case identifier contributed at most one paired value to each
Baseline-vs-preictal contrast.

Post-hoc comparisons were interpreted cautiously because a large number of
contrasts were tested across graph metrics, connectivity methods, and frequency
bands. In the final interpretation, post-hoc results were considered most
meaningful when the corresponding omnibus temporal-window effect was
significant after Benjamini--Hochberg FDR correction. This rule reduced the risk
of over-interpreting isolated pairwise differences that were not supported by
an overall window effect.

\section{Preictal Trend Analysis}
\label{sec:preictal-trend-analysis}

In addition to comparing each preictal window with Baseline, the statistical
pipeline tested whether graph metrics changed progressively during the
preictal period itself. This analysis used only the three preictal windows T0,
T1, and T2. Baseline was not included because the purpose of this test was to
evaluate directional change as seizure onset approached, rather than to compare
the preictal period with the interictal reference condition.

The three preictal windows were represented by their midpoint times relative to
seizure onset. T0 was assigned a time of \(-7.5\) minutes, T1 a time of
\(-4.5\) minutes, and T2 a time of \(-1.5\) minutes. These values correspond to
the centers of the three consecutive three-minute preictal windows. A positive
slope therefore indicates that the graph metric increased as the recording
approached seizure onset, whereas a negative slope indicates that the metric
decreased across the preictal period.

For each graph metric, connectivity method, and frequency band combination, the
primary preictal trend analysis used a linear mixed-effects model when
available. The model can be written as

\[
    y_{s,t}
    =
    \beta_0
    +
    \beta_1 t
    +
    u_{\mathrm{subject}}
    +
    u_{\mathrm{seizure}}
    +
    \varepsilon_{s,t},
\]

where $y_{s,t}$ is the graph metric value for seizure $s$ at preictal time
$t$, $\beta_0$ is the intercept, $\beta_1$ is the fixed slope across preictal
time, $u_{\mathrm{subject}}$ accounts for subject/case-level grouping,
$u_{\mathrm{seizure}}$ accounts for repeated measurements within seizure, and
$\varepsilon_{s,t}$ is the residual error term.

The main parameter of interest was the slope coefficient $\beta_1$. This
coefficient describes the average change in the graph metric per minute during
the preictal interval. Because time was coded in minutes relative to seizure
onset, the sign of the slope must be interpreted with the direction of time in
mind. A positive slope means that the metric increased from T0 toward T2, while
a negative slope means that the metric decreased as seizure onset approached.

Only seizures with complete T0, T1, and T2 values were included in the
seizure-level preictal trend analysis. This ensured that the trend was
estimated from within-seizure trajectories across all three preictal windows.
The same complete-case logic was applied to each metric--method--band
combination separately.

When the mixed-effects model could not be used, the pipeline applied a fallback
trend analysis. In the fallback approach, a linear slope was first estimated
separately for each seizure across T0, T1, and T2. These seizure-level slopes
were then tested against zero using a one-sample $t$-test. This preserved the
interpretation of the trend as a within-seizure temporal change, although it
did not explicitly model subject/case-level grouping.

The trend analysis was intended to complement the omnibus and post-hoc tests.
The omnibus test assessed whether any difference existed across Baseline, T0,
T1, and T2. The post-hoc tests identified which preictal windows differed from
Baseline. The trend test asked a different question: whether the metric changed
systematically within the preictal period itself. Therefore, a metric could
show a Baseline-vs-preictal difference without a monotonic T0--T2 trend, or it
could show a preictal trend without a strong Baseline contrast.

For each trend test, the statistical output recorded the number of complete
preictal seizures, the number of contributing subject/case identifiers, the
model type, the estimated slope per minute, the test statistic, the
uncorrected $p$-value, the FDR-corrected $p$-value, and any fallback or skip
reason. As with the other statistical outputs, trend results were interpreted
after correction for multiple comparisons.

\section{Multiple-Comparison Correction}
\label{sec:multiple-comparison-correction}

Because graph metrics were tested separately across connectivity methods and
frequency bands, the analysis involved many statistical tests. To reduce the
risk of false-positive interpretation, Benjamini--Hochberg false discovery rate
(FDR) correction was applied \cite{Benjamini1995}. FDR correction controls the
expected proportion of false discoveries among results declared significant,
which is appropriate for an exploratory analysis with many related tests.

Correction was applied separately within each output family: seizure-level
omnibus tests, seizure-level post-hoc tests, seizure-level preictal trend
tests, subject-average omnibus tests, and subject-average post-hoc tests. Thus,
correction was performed across tests belonging to the same statistical family,
rather than separately for each subject/case identifier or each individual graph
metric.

For post-hoc analyses, paired $t$-test $p$-values and Wilcoxon signed-rank
$p$-values were corrected separately. A corrected significance threshold of
$\alpha = 0.05$ was used. Uncorrected $p$-values were retained for
transparency, but interpretation focused on FDR-corrected values. Planned
post-hoc comparisons were interpreted primarily when the corresponding omnibus
window-effect test also survived FDR correction.

\section{Validation and Output of the Statistical Pipeline}
\label{sec:statistical-pipeline-validation-output}

The statistical pipeline transformed the graph feature table into a long-format
dataset suitable for repeated-measures testing. The input contained
\num{11376} graph-level rows and \num{8} analyzed graph metrics, resulting in
\num{91008} long-format observations. Each observation retained its
subject/case identifier, seizure identifier, temporal window, frequency band,
connectivity method, graph metric, and value.

The final dataset included \num{24} subject/case identifiers and \num{159}
usable seizures. Of these, \num{155} seizures had complete values across
Baseline, T0, T1, and T2 and were used for Baseline-inclusive seizure-level
omnibus testing. All \num{159} seizures had complete T0, T1, and T2 values and
were available for the preictal trend analysis.

The analysis covered \num{144} metric--method--band combinations
(\num{8} graph metrics $\times$ \num{3} connectivity methods $\times$
\num{6} frequency bands). At the seizure level, this produced \num{144}
omnibus tests, \num{432} planned Baseline-vs-preictal post-hoc contrasts, and
\num{144} preictal trend tests. The subject-average sensitivity analysis
produced \num{13824} rows, \num{112} omnibus tests with valid model output, and
\num{432} post-hoc contrasts.

The output tables included descriptive summaries, seizure-level omnibus tests,
seizure-level post-hoc comparisons, seizure-level preictal trends,
subject-average summaries, subject-average omnibus tests, and subject-average
post-hoc comparisons. Inferential outputs retained the model type, test
statistic, uncorrected and FDR-corrected $p$-values, effect-size information
where applicable, and diagnostic fields such as fallback or skip reasons.

These validation outputs confirmed that the dataset structure was preserved
from graph extraction to statistical testing and provided the basis for the
results chapter.